% Generated draft: requires mathematical and source review.
\documentclass[11pt]{article}
\usepackage[margin=25mm]{geometry}
\usepackage{fontspec}
\setmainfont{Latin Modern Roman}
\usepackage{amsmath,amssymb,mathtools,unicode-math}
\setmathfont{Latin Modern Math}
\usepackage{enumitem,needspace,xurl,hyperref}
\setlength{\parindent}{0pt}
\setlength{\parskip}{5pt}
\setlength{\emergencystretch}{3em}
\newcommand{\N}{\mathbb{N}}
\newcommand{\Z}{\mathbb{Z}}
\newcommand{\Q}{\mathbb{Q}}
\newcommand{\R}{\mathbb{R}}
\newcommand{\C}{\mathbb{C}}
\newcommand{\F}{\mathbb{F}}
\newcommand{\A}{\mathbb{A}}
\newcommand{\PP}{\mathbb{P}}
\newcommand{\E}{\mathbb{E}}
\newcommand{\eps}{\varepsilon}
\newcommand{\dd}{\,\mathrm{d}}
\begin{document}
\section*{UnsolvedMath ID: 30000067}
\textbf{Source problem code:} catalogue entry
\subsection{Definitions and mathematical statement}
The source asks whether minimizers $(u_\varepsilon,\rho_\varepsilon)$ of a surfactant energy $G_\varepsilon$ converge to a macroscopic state as $\varepsilon\to0$, and whether the limit minimizes a relaxed macroscopic energy. The functional is not given. Neither the spatial domain, admissible classes, interpretation and constraints on $u,\rho$, boundary data nor convergence topology are specified. No phase-field or surfactant coupling formula is invented.

To state the requested conclusions abstractly, let $z_\varepsilon=(u_\varepsilon,\rho_\varepsilon)$ belong to an admissible state space with a specified topology $\tau$, and suppose
\[
 z_\varepsilon\in\operatorname*{argmin}G_\varepsilon.
\]
The compactness target would be that every suitable sequence $\varepsilon_j\downarrow0$ of such minimizers has a $\tau$-convergent subsequence. The identification target would be that each resulting cluster point $z$ satisfies
\[
 z\in\operatorname*{argmin}G_{\mathrm{rel}},
\]
for an explicitly defined relaxed macroscopic functional and admissible constraints. These express possible precise meanings of the question, not conclusions established for the missing model. Subsequence convergence and convergence of the entire family differ; the latter generally needs additional selection or uniqueness information.

A standard variational route would establish equicoercivity, an appropriate $\Gamma$-limit, matching boundary/mass constraints and finite comparison energies before transferring global minimizers to the limit. None of those hypotheses or a candidate $G_{\mathrm{rel}}$ is supplied. A relaxation is relative to a topology and initial functional; it is not automatically identical to an asserted $\Gamma$-limit. Weak function-space convergence and measure convergence allowing surfactant concentration could also give different state spaces and limits. Local minimizers would require a separate formulation.

The supplied literature triage checked 21 August 2026 explicitly says the energy is insufficiently specified to verify this as a distinct open problem using available relative-isoperimetric or $\Gamma$-convergence results. The open label therefore does not certify a complete well-posed model or a new unresolved theorem. The draft preserves the compactness and limiting-minimization questions with their missing data rather than claiming a general convergence result.
\subsection{Short English statement}
The requested macroscopic-limit statement cannot be made model-specific without the energy, constraints and topology. Compactness, identification of limit minimizers and full-family convergence are separate goals. The dated review itself says the extracted formulation is insufficient to verify a distinct open problem.
\subsection{Sources}
\begin{itemize}
\item Local UnsolvedMath registry, permanent ID 30000067, source code catalogue entry, supplied Oberwolfach Reports record; updated 16 September 2026.
\item Supplied citation: \emph{Calculus of Variations (2005)}, Oberwolfach Reports, DOI 10.4171/owr/2004/29, pp. 1541--1586: \url{https://doi.org/10.4171/owr/2004/29}. The source separately lists proposed year 2004.
\item The only supplied literature reference is the original report at the same DOI. No energy formula, model-specific theorem or additional authoritative citation is supplied.
\item Literature triage checked 21 August 2026, derived from the supplied OWR curation; independently verify time-sensitive status and cited claims. It explicitly records insufficient model data.
\end{itemize}
\end{document}
